% Part: methods
% Chapter: proofs
% Section: using-definitions

\documentclass[../../../include/open-logic-section]{subfiles}

\begin{document}

\olfileid[pt]{mod}{prf}{def}

\olsection{Usando Definições}

Mencionamos que você deve estar familiarizado com todas as definições que possam
ser usadas na prova, e que saiba aplicá-las corretamente. Esse é um
ponto realmente importante, e vale a pena examiná-lo em um pouco mais de
detalhe. As definições são usadas para abreviar propriedades e relações, de modo que
possamos falar delas de forma mais sucinta. A abreviação introduzida é
chamada de \emph{definiendum}, e aquilo que ela abrevia é o
\emph{definiens}.  Em provas, com frequência temos que voltar a como o
definiendum foi introduzido, porque temos que explorar a estrutura
lógica do definiens (a versão longa da qual o termo definido
é a abreviação) para concluir nossa prova.  Ao desdobrar as
definições, você garante que está chegando ao cerne de onde
a ação lógica acontece.

Começaremos com um exemplo. Suponha que você queira provar o seguinte:

\begin{prop}
Para quaisquer conjuntos $A$ e $B$, $A \cup B = B \cup A$.
\end{prop}

Para sequer iniciar a prova, precisamos saber o que significa que
dois conjuntos sejam idênticos; isto é, precisamos saber o que o ``$=$'' naquela
equação significa para conjuntos.  Os conjuntos são definidos como idênticos
sempre que têm os mesmos !!{element}s.  Assim, a definição que temos
que desdobrar é:

\begin{defn}
Os conjuntos $A$ e $B$ são \emph{idênticos}, $A = B$, sse todo !!{element}
de~$A$ é !!a{element} de~$B$, e vice-versa.
\end{defn}

Esta definição usa $A$ e~$B$ como marcadores de posição para conjuntos
arbitrários. O que ela define---o \emph{definiendum}---é a expressão
``$A = B$'', dando a condição sob a qual $A = B$ é verdadeira.  Essa
condição---``todo !!{element} de~$A$ é !!a{element} de~$B$, e
vice-versa''---é o \emph{definiens}.\footnote{Neste caso
  particular---e de modo bastante confuso!---quando $A = B$, os conjuntos $A$ e $B$
  são apenas um único e mesmo conjunto, ainda que usemos letras diferentes
  para ele do lado esquerdo e do lado direito.  Mas os modos pelos quais esse
  conjunto é identificado podem ser diferentes, e isso torna a definição
  não trivial.} A definição especifica que $A = B$ é verdadeira se, e
somente se (abreviamos isso por ``sse'') a condição vale.

Quando você aplica a definição, tem que fazer corresponder o $A$ e o $B$ da
definição ao caso de que está tratando.  No nosso caso, isso significa
que, para que $A \cup B = B \cup A$ seja verdadeira, cada $z \in A \cup
B$ também deve estar em $B \cup A$, e vice-versa.  A expressão $A \cup
B$ na proposição faz o papel de~$A$ na definição, e $B
\cup A$, o de~$B$. Como $A$ e $B$ são usados tanto na definição
quanto no enunciado da proposição que estamos provando, mas em
usos diferentes, você precisa ter cuidado para não confundir
os dois.  Por exemplo, seria um erro pensar que você poderia
provar a proposição mostrando que todo !!{element} de~$A$ é
!!a{element} de~$B$, e vice-versa---isso mostraria que $A = B$,
não que $A \cup B = B \cup A$. (Além disso, como $A$ e $B$ podem ser quaisquer
dois conjuntos, você não chegaria muito longe, pois, se nada for suposto sobre
$A$ e~$B$, eles bem podem ser conjuntos diferentes.)

Dentro da prova, estamos lidando com noções da teoria dos conjuntos, como
união, e portanto também precisamos conhecer os significados do símbolo $\cup$ a fim
de entender como a prova deve prosseguir. E, às vezes,
desdobrar a definição dá origem a outras definições a
desdobrar. Por exemplo, $A \cup B$ é definido como $\Setabs{z}{z \in A
  \text{ ou } z \in B}$. Então, se você quer provar que $x \in A \cup
B$, desdobrar a definição de $\cup$ lhe diz que você tem que
provar $x \in \Setabs{z}{z \in A \text{ ou } z \in B}$.  Agora você também
tem que lembrar que $x \in \Setabs{z}{\dots z\dots}$ sse $\dots
x\dots$.  Assim, desdobrando ainda mais a definição da
notação $\Setabs{z}{\dots z \dots}$, o que você tem que mostrar é: $x \in
A$ ou $x \in B$. Assim, ``todo !!{element} de $A \cup B$ é também
!!a{element} de $B \cup A$'' significa, na verdade: ``para todo $x$, se $x \in
A$ ou $x \in B$, então $x \in B$ ou $x \in A$.''  Se desdobramos
completamente as definições na proposição, vemos que o que temos que
mostrar é isto:

\begin{prop}
Para quaisquer conjuntos $A$ e $B$: (a) para todo $x$, se $x \in A$ ou $x \in
B$, então $x \in B$ ou $x \in A$, e (b) para todo $x$, se $x \in B$
ou $x \in A$, então $x \in A$ ou $x \in B$.
\end{prop}

O que importa é que desdobrar definições é uma parte necessária da
construção de uma prova. Fazê-lo corretamente é às vezes difícil: você
deve ter cuidado para distinguir e fazer corresponder as variáveis da
definição e os termos na afirmação que está provando.  Para ter
sucesso, você deve saber o que a questão está pedindo e o que todos os
termos usados na questão significam---muitas vezes você precisará desdobrar mais
de uma definição.  Em provas simples como as que seguem, a
solução decorre quase imediatamente das próprias
definições. É claro que nem sempre será assim tão simples.

\begin{prob}
Suponha que lhe peçam para provar que $A \cap B \neq \emptyset$. Desdobre
todas as definições que ocorrem aqui, isto é, reenuncie isto de um modo que
não mencione ``$\cap$'', ``='', ou ``$\emptyset$''.
\end{prob}

\end{document}
